\section{Аналитическая часть}

Проблемы эмулятора:
\begin{itemize}[leftmargin=1.6\parindent]
	\item[---] управление кэшем транслированного кода;
	\item[---] выделение регистров;
	\item[---] оптимизация условных блоков;
	\item[---] direct block chaining????? lmao;
	\item[---] управление памятью;
	\item[---] поддержка самоизменяемого кода;
	\item[---] поддержка исключений;
	\item[---] поддержка аппаратных прерываний.
\end{itemize}

\subsection{QEMU}

кему это круто написать что это такое

\subsubsection{Сворачивание и оптимизация тривиальных выражений}

В qemu определены свои опкоды, каждый из них может по своему сворачиваются. TCG ops (специальные операции TCG) оптимизируются и затем выполняются на хосте, таким образом портировать TCG становится проще. Операции поддерживаются над 32 и 64 битными целыми числами, указатели определены как алиас над соответствующим целым числом.

Пример операции TCG:\\
\texttt{add\_i32 t0, t1, t2  (t0 <- t1 + t2)}

Одной из оптимизаций является игнорирование бессмысленных операций, например операция:\\
\texttt{and\_i32 t0, t0, \$0xffffffff\\}
выполняться не будет

Подавляются неиспользуемые перемещения данных, например из трех операций:\\
\texttt{add\_i32 t0, t1, t2\\
add\_i32 t0, t0, \$1\\
mov\_i32 t0, \$1\\}
выполнится только последняя.

(оптимизации над инструкциями реализованы в файле tcg/optimize.c)

\subsubsection{Состояние процессора и блоки трансляции}

Блоком трансляции называется часть кода до момента изменения состояния процессора которое нельзя выяснить на этапе трансляции (например, некоторое ветвление).

Каждый регистр проверяется на используемость в определенном блоке трансляции, если регистр не используется в блоке трансляции связанные с ним операции оптимизируются, так как значения этих регистров (и связанное с ними состояние процессора) за блок не могут измениться. Если состояние виртуального процессора меняется, такой блок не будет выполняться пока состояние процессора не будет соответствовать необходимому для блока (например, другой уровень привилегий).

Пример: если на x86 регистры SS, DS и SS содержат в себе 0, транслятор не будет генерировать для них смещение.

\subsubsection{Блок чейнинг??}

После выполнения блока трансляции QEMU ищет следующий блок, для этого используется PC (эмулируемый program counter) и другая информация о статусе процессора (например регистр CS). (реализовано в функции tb\_lookup в файле accel/tcg/cpu-exec.c)

Сначала блок ищется в кэше трансляции, если он там не найден он достается из хэш-таблицы и добавляется в кэш.

Для этого нужно выйти из цикла выполнения блока, пройти через эпилог процедуры, найти следующий блок, запустить его, пройдя через пролог процедуры. В качестве оптимизации предлагается связывать несколько блоков трансляции напрямую.

Для этого в конце блока вызывается tcg\_gen\_lookup\_and\_goto\_ptr(), он в свою очередь вызывает helper\_lookup\_tb\_ptr который ищет нужный блок и генерирует инструкцию goto\_ptr, которая либо продолжит управление в нужном блоке, либо вернется в основной цикл. (пример вызова для i386 в файле target/i386/tcg/translate.c в функции do\_gen\_eob\_worker)

Еще одной оптимизацией является оптимизация ветвления. Если ветвление происходит напрямую, в пределах одной страницы QEMU выполняет поиск блока трансляции на который будет произведено ветвление, а затем сохраняет его адрес в транслированном коде. Таким образом при следующем выполнении этого блока нет необходимости в поиске следующего блока.

\subsubsection{Поддержка самомодифицирующегося кода}

Самомодифицирующийся код на x86 представляет особую проблему, так как нет механизма оповещения о изменении кода. При запуске в режиме пользователя QEMU помечает все страницы с транслированным кодом как защищенные от записи, при попытке записи в них поднимается сигнал SEGV, допускается запись, обнуляются все транслированные страницы и связь блоков. При запуске эмуляции системы программный MMU выполняет защиту от записи.
Для эффективного выполнения этих операций хранится связанный список всех блоков трансляции и связей блоков.

\subsubsection{Эмуляция MMU}

Каждый доступ к памяти проходит через MMU, в нем адреса преобразуются из виртуальных к реальным. Для ускорения трансляции адреса хранится TLB-кэш.

Все блоки трансляции в кэше индексируются при помощи их физического адреса, таким образом нет необходимости очищать кэш при смене страницы.

\subsection{bruh}

Список:

\begin{itemize}[leftmargin=1.6\parindent]
	\item[---] первое;
	\item[---] второе;
	\item[---] пятое;
	\item[---] десятое.
\end{itemize}

Формула:

\begin{equation}
c^2 = a^2 + b^2
\end{equation}

Ссылаемся на рисунок \ref{fig:a1}. Информация из источника \cite{golang}.

\begin{figure}[hbtp]
	\centering
	\includegraphics[width=\textwidth]{img/golang.png}
	\caption{Пример рисунка}
	\label{fig:a1}
\end{figure}

\begin{code}
	\captionof{listing}{Пример кода}
	\label{code:1}
	\inputminted
	[
	frame=single,
	framerule=0.5pt,
	framesep=20pt,
	fontsize=\small,
	tabsize=4,
	linenos,
	numbersep=5pt,
	xleftmargin=10pt,
	]
	{text}
	{code/main.go}
\end{code}

\pagebreak